\chapter{Symplectic Manifolds}

\section{Symplectic Manifolds}

\begin{definition}
    By a \textbf{symplectic manifold} we mean an even-dimensional differentiable $\left(C^{\infty}\right)$ manifold $M^{2 n}$ together with a global 2-form $\Omega$ which is closed and of maximal rank, i.e., $d \Omega=0, \Omega^n \neq 0$. 
    
    By a \textbf{symplectomorphism} $f:\left(M_1, \Omega_1\right) \longrightarrow\left(M_2, \Omega_2\right)$ we mean a diffeomorphism $f: M_1 \longrightarrow M_2$ such that 
    $$f^* \Omega_2=\Omega_1.$$
\end{definition}

\begin{definition}
    Let $\iota: L \longrightarrow M^{2 n}$ be an immersion into a symplectic manifold $\left(M^{2 n}, \Omega\right)$. We say that $L$ is a \textbf{Lagrangian submanifold} if the dimension of $L$ is $n$ and $\iota^* \Omega=0$.
\end{definition}

\begin{theorem}[Weistein]
    If $L$ is a Lagrangian submanifold of a symplectic manifold $(M, \Omega)$, then there exists a neighborhood of $L$ in $M$ that is symplectomorphic to a neighborhood of the zero section in $T^{*}L$.
\end{theorem}

\begin{theorem}
    Let $X$ be a symplectic vector field on $(M, \Omega), g$ an associated metric, and $J$ the corresponding almost complex structure. Then $X^i=J^{i k} \theta_k$ for some closed 1-form $\theta$. Conversely, given a closed 1-form $\theta, X^i=J^{i k} \theta_k$ defines a symplectic vector field.
\end{theorem}
\begin{corollary}
    For $f \in C^{\infty}(M), J \nabla f$ is symplectic, where $\nabla f$ is the gradient of $f$. Conversely, given $X$ symplectic, $X$ is locally $J \nabla f$.
\end{corollary}


\begin{theorem}[Hatakeyama, 1966]
    The group of symplectomorphisms acts transitively on a compact symplectic manifold $(M, \Omega)$.
\end{theorem}


\subsection{Examples}
\begin{example}
    Let $M$ be a differentiable manifold. Then its cotangent bundle $T^* M$ has a natural symplectic structure. For $z \in T^* M$ and $V$ in the tangent space of $T^* M$ at $z, T_z T^* M$, define a 1 -form $\beta$, often called the \textbf{Liouville form}, by $\beta(V)_z=z\left(\pi_* V\right)$, where $\pi: T^* M \longrightarrow M$ is the projection map. If $x^1, \ldots, x^n$ are local coordinates on $M$, then $q^i=x^i \circ \pi$ together with fiber coordinates $p^1, \ldots, p^n$ give local coordinates on $T^* M$. In these coordinates $\beta$ has the local expression $\sum_{i=1}^n p^i d q^i$. The natural symplectic structure on $T^* M$ is given by $\Omega=-d \beta$.
\end{example}



\section{Contact manifolds}
\begin{definition}
    By a \textbf{contact manifold} we mean a $C^{\infty}$ manifold $M^{2 n+1}$ together with a 1 -form $\eta$ such that 
    \begin{equation}
        \eta \wedge(d \eta)^n \neq 0. 
    \end{equation}
\end{definition}
\begin{remark}
   1)  In particular, $\eta \wedge(d \eta)^n \neq 0$ is a volume element on $M$, so that a contact manifold is orientable. 
    
   2) we have a global vector field $\xi$ satisfying
    \begin{equation}
        d \eta(\xi, X)=0, \quad \eta(\xi)=1 .
    \end{equation}
$\xi$ is called the \textbf{characteristic vector field} or \textbf{Reeb vector field } of the contact structure $\eta$. 

3) By formula $\left.\left.L_{\xi}=d \circ(\xi\lrcorner\right)+(\xi\lrcorner\right) \circ d$, we have immediately that
\begin{equation}
    L_{\xi} \eta=0, \quad L_{\xi} d \eta=0 .
\end{equation}
\end{remark}

\begin{definition}
    We denote by $\mathcal{D}$ the \textbf{contact distribution} or subbundle defined by the subspaces $\mathcal{D}_m=\left\{X \in T_m M: \eta(X)=0\right\}$. 
\end{definition}
\begin{remark}
    Roughly speaking, the meaning of the contact condition, $\eta \wedge(d \eta)^n \neq 0$, is that the contact subbundle is as far from being integrable as possible; in particular, $\mathcal{D}$ rotates as one moves around on the manifold. For a subbundle defined by a 1 -form $\eta$ to be integrable it is necessary and sufficient that 
    \begin{equation}
        \eta \wedge(d \eta) \equiv 0. 
    \end{equation}
\end{remark}

\begin{definition}
    A contact structure is \textbf{regular} if $\xi$ is \textbf{regular} as a vector field, that is, every point of the manifold has a neighborhood such that any integral curve of the vector field passing through the neighborhood passes through only once. 
\end{definition}
\begin{remark}
    Two well-known examples of nonregular vector fields on surfaces are the irrational flow on a torus and the flow around a Möbius band.
\end{remark}

\begin{theorem}[\textbf{Darboux coordinate system}]
    About each point of a contact manifold $\left(M^{2 n+1}, \eta\right)$ there exist local coordinates $\left(x^1, \ldots, x^n, y^1, \ldots, y^n, z\right)$ with respect to which
    \begin{equation}
        \eta=d z-\sum_{i=1}^n y^i d x^i.
    \end{equation}
\end{theorem}
\begin{remark}
    There are many similar results in this spirit.
\end{remark}

\begin{definition}
    A diffeomorphism $f$ of $M^{2 n+1}$ or between open subsets of $\mathbb{R}^{2 n+1}$ with the contact structure of the Darboux form of the above theorem is called a \textbf{contact transformation} if 
    \begin{equation}
        f^* \eta=\tau \eta
    \end{equation}
    for some nonvanishing function $\tau$ on the domain of $f$. If $\tau \equiv 1, f$ is called a \textbf{strict contact transformation}.
\end{definition}

\begin{definition}
    There is also the notion of a \textbf{contact structure in the wider sense,} often called simply a "contact structure" by many authors, which can be defined in a number of ways. For example, a contact manifold in the wider sense is a manifold with a differentiable structure modeled on the pseudogroup of contact transformations on $\mathbb{R}^{2 n+1}$.
\end{definition}

\begin{theorem}
 Let $M^{2 n+1}$ be a contact manifold in the wider sense. If $n$ is odd, then $M^{2 n+1}$ is orientable.
\end{theorem}
\begin{theorem}
     Let $M^{2 n+1}$ be a contact manifold in the wider sense. If $n$ is even and $M^{2 n+1}$ is orientable, then $M^{2 n+1}$ is a contact manifold.
\end{theorem}

\begin{theorem}
    Let $M^{2 n+1}$be a contact manifold in the wider sense which is not a contact manifold in the restricted sense. Then its 2-fold covering manifold is a contact manifold in the restricted sense.
\end{theorem}

\begin{theorem}[\textbf{Boothby-Wang fibration}]
     Let $\left(M^{2 n+1}, \eta^{\prime}\right)$ be a compact regular contact manifold. Then there exists a contact form $\eta=\tau \eta^{\prime}$ for some nonvanishing function $\tau$ whose characteristic vector field $\xi$ generates a free effective $S^1$ action on $M^{2 n+1}$. Moreover, $M^{2 n+1}$ is the bundle space of a principal circle bundle $\pi: M^{2 n+1} \longrightarrow M^{2 n}$ over a symplectic manifold $M^{2 n}$ whose symplectic form $\Omega$ determines an integral cocycle on $M^{2 n}$ and $\eta$ is a connection form on the bundle with curvature form $d \eta=\pi^* \Omega$.
\end{theorem}

\begin{definition}[Other notions]
    \textbf{periodic point, planar, almost regular/quasi-regular, R-contact manifold}, etc..
\end{definition}

\subsection{Examples}
\begin{example}
    The cotangent sphere bundle and the tangent sphere bundle of a Riemannian manifold are contact manifolds.
\end{example}

\begin{example}
    Let $M$ be an $n$-dimensional manifold and $T^* M$ its cotangent bundle. As in the previous example we can define a 1 -form $\beta$ by the local expression $\beta=\sum_{i=1}^n p^i d q^i$. Let $M^{2 n+1}=T^* M \times \mathbb{R}, t$ the coordinate on $\mathbb{R}$, and $\mu: M^{2 n+1} \longrightarrow T^* M$ the projection to the first factor. Then $\eta=d t-\mu^* \beta$ is a contact form on $M^{2 n+1}$.
\end{example}

\section{Associated Metrics}
\begin{definition}
    An \textbf{almost complex structure} is a tensor field $J$ of type $(1,1)$ such that $J^2=-I$. A \textbf{Hermitian metric} on an almost complex manifold $(M, J)$ is a Riemannian metric that is invariant by $J$, i.e.,
    \begin{equation}
        g(J X, J Y)=g(X, Y).
    \end{equation}
    \end{definition}

    \begin{remark}
        Note also that every almost complex manifold admits a Hermitian metric, for if $k$ is any Riemannian metric, then $g$ defined by
        \begin{equation}
            g(X, Y)=k(X, Y)+k(J X, J Y)
        \end{equation}
is Hermitian.
    \end{remark}
    
    \begin{definition}
    Noting that $J$ is negative self-adjoint with respect to $g$, i.e., $g(X, J Y)=$ $-g(J X, Y)$, therefore 
    \begin{equation}
         \Omega(X, Y)=g(X, J Y)
    \end{equation}
    defines a 2 -form called the\textbf{ fundamental 2-form} of the \textbf{almost Hermitian structure} $(M, J, g)$. If 
    \begin{equation}
            d \Omega= 0,
    \end{equation}
 the structure is \textbf{almost Kähler}. If $M$ is a complex manifold and $J$ the corresponding almost complex structure, we say that $(M, J, g)$ is a \textbf{Hermitian manifold}.
\end{definition}
\begin{definition}
    Given $(M, J, g)$ we can construct a particular local orthonormal basis as follows. Let $\mathcal{U}$ be a coordinate neighborhood on $M$ and $X_1$ any unit vector field on $\mathcal{U}$. Let $X_{1^*}=J X_1$. Now choose a unit vector field $X_2$ orthogonal to both $X_1$ and $X_{1^*}$. Then $J X_2$ is also orthogonal to $X_1$ and $X_{1^*}$. Continuing in this manner we have a local orthonormal basis of the form $\left\{X_1, \ldots, X_n, J X_1, \ldots, J X_n\right\}$. Such a basis is called a \textbf{$J$-basis.} 
\end{definition}

\begin{remark}
    An almost complex (almost Hermitian) structure on $M$ can be thought of as a reduction of the structural group to $U (n)$.
\end{remark}

\begin{remark}
    The structural group of a symplectic manifold is reducible to $U(n)$ and that of a contact manifold to $U(n) \times 1$ (Chern [1953]). 
\end{remark}
\begin{definition}
        For an odd-dimensional manifold $M^{2 n+1}$, J. Gray [1959] defined \textbf{an almost contact structure} as a reduction of the structural group to $U(n) \times 1$. In terms of structure tensors we say that $M^{2 n+1}$ has an almost contact structure or sometimes \textbf{$(\phi, \xi, \eta)$-structure} if it admits a tensor field $\phi$ of type $(1,1)$, a vector field $\xi$, and a 1 -form $\eta$ satisfying
    \begin{equation}
        \phi^2=-I+\eta \otimes \xi, \quad \eta(\xi)=1 .
    \end{equation}
\end{definition}
\begin{theorem}
    Suppose $M^{2 n+1}$ has a $(\phi, \xi, \eta)$-structure. Then 
    \begin{equation}
        \phi \xi=0, \quad \eta \circ \phi=0.
    \end{equation}
 Moreover, the endomorphism $\phi$ has rank $2 n$.
\end{theorem}

\begin{definition}
    If a manifold $M^{2 n+1}$ with a $(\phi, \xi, \eta)$-structure admits a Riemannian metric $g$ such that
    \begin{equation}\label{eq12}
        g(\phi X, \phi Y)=g(X, Y)-\eta(X) \eta(Y),
    \end{equation}
we say that $M^{2 n+1}$ has an \textbf{almost contact metric structure} and $g$ is called a \textbf{compatible metric}. 
\end{definition}
\begin{remark}
    For any $X\in \mathcal{D}$, let $Y=\xi$ in $\eqref{eq12}$, we get 
    $g(\phi(X),\xi)=0$, thus
    \begin{equation}
        \phi (\mathcal{D}) \subset \mathcal{D}.
    \end{equation}
\end{remark}

\begin{remark}
        As in the almost Hermitian case, we define a 2 -form, called the fundamental 2-form of the almost contact metric structure, by 
        \begin{equation}
            \Phi(X, Y)=g(X, \phi Y). 
        \end{equation}
\end{remark}
\begin{remark}
    Also as in the case of an almost complex structure, the existence of the compatible metric is easy. For if $k^{\prime}$ is any metric, first set $k(X, Y)=k^{\prime}\left(\phi^2 X, \phi^2 Y\right)+\eta(X) \eta(Y)$; then $\eta(X)=k(X, \xi)$. Now define $g$ by
    \begin{equation}
        g(X, Y)=\frac{1}{2}(k(X, Y)+k(\phi X, \phi Y)+\eta(X) \eta(Y)).
    \end{equation}
\end{remark}


\begin{remark}
    Setting $Y=\xi$ we have immediately that
\begin{equation}
    \eta(X)=g(X, \xi).
\end{equation}
\end{remark}
\begin{remark}
    We can obtain a local orthonormal basis $\left\{X_i, X_{i^*}=\phi X_i, \xi\right\}$, called a \textbf{$\phi$-basis}.
\end{remark}
\begin{remark}
    Given a manifold $M^{2 n+1}$ with a $(\phi, \xi, \eta)$-structure, let $g$ be a compatible metric, then proceeding as in the almost complex case, we see that the structural group of $M^{2 n+1}$ is reducible to $U(n) \times 1$. 
\end{remark}

\begin{theorem}[\textbf{Existence of associated metrics 1}]
    Let $\left(M^{2 n}, \Omega\right)$ be a symplectic manifold. Then there exist a Riemannian metric $g$ and an almost complex structure $J$ such that
$$
g(X, J Y)=\Omega(X, Y) .
$$
\end{theorem}

\begin{definition}
    In particular, given a symplectic manifold $\left(M^{2 n}, \Omega\right)$, we say that a Riemannian metric $g$ is an \textbf{associated metric} if there exists an almost complex structure $J$ such that 
    $$g(X, J Y)=\Omega(X, Y).$$
\end{definition}

\begin{theorem}[\textbf{Existence of associated metrics 2}]
    Let $\left(M^{2 n+1}, \eta\right)$ be a contact manifold and $\xi$ its characteristic vector field. Then there exists an almost contact metric structure such that 
    $$g(X, \phi Y)=d \eta(X, Y).$$
\end{theorem}

\begin{definition}
    As in the symplectic case, given a contact manifold $\left(M^{2 n+1}, \eta\right)$, we say that a Riemannian metric $g$ is an \textbf{associated metric} if there exists an almost contact metric structure such that 
   $$g(X, \phi Y)=d \eta(X, Y).$$
    In this case we also speak of a\textbf{ contact metric structure}; other authors often use the phrase contact Riemannian structure. 
\end{definition}
\begin{definition}[Alternative definition]
    Working strictly with structure tensors, one may avoid the polarization process by defining an associated metric as follows: Given a contact manifold $\left(M^{2 n+1}, \eta\right)$ with characteristic vector field $\xi$, a Riemannian metric $g$ is an associated metric if, first of all,
$$
\eta(X)=g(X, \xi)
$$
and secondly, there exists a tensor field $\phi$ of type $(1,1)$ such that
$$
\phi^2=-I+\eta \otimes \xi, \quad d \eta(X, Y)=g(X, \phi Y) .
$$
\end{definition}
\begin{proof}
    Only need to verify 
    \begin{equation}
        g(X,\phi Y)+g(\phi  X,Y)=0, X,Y \in \mathcal{D}.
    \end{equation}
    But it follows immediately by symmetry of $g$ and antisymmetry of $d\eta$.
\end{proof}
\begin{remark}
    Caution that it is possible to have a contact manifold $\left(M^{2 n+1}, \eta\right)$ with characteristic vector field $\xi$ and an almost contact metric structure $(\phi, \xi, \eta, g)$, same $\xi$ and $\eta$, without $g(X, \phi Y)=d \eta(X, Y)$.
\end{remark}
\begin{remark}
    On the other hand, it is possible for a Riemannian metric g to be an associated metric for more than one symplectic structure. For example, on a hyper-Kähler manifold one has, by definition, three independent global Kähler structures $\left(J_a, g\right), a=1,2,3$, satisfying $J_1 J_2+J_2 J_1=0$, $J_3=J_1 J_2$. Thus the three fundamental 2 -forms give three symplectic structures with $g$ an associated metric for each of them.
\end{remark}

\begin{theorem}
     On a contact metric manifold the integral curves of $\xi$ are geodesics.
\end{theorem}
\begin{proof}
    For a contact metric structure we have
$$
0=\left(L_{\xi} \eta\right)(X)=\xi g(X, \xi)-g\left(\nabla_{\xi} X-\nabla_X \xi, \xi\right)=g\left(X, \nabla_{\xi} \xi\right),
$$
so the integral curves of $\xi$ are geodesics.
\end{proof}


\begin{remark}
    The set $\mathcal{A}$ of all associated metrics forms a infinite-dimensional sapce.
\end{remark}

\begin{theorem}
     Let $\left(M^{2 n}, \Omega\right)$ be a symplectic manifold (resp. $\left(M^{2 n+1}, \eta\right)$ a contact manifold) and $g$ an associated metric. Then
     \begin{equation}
     d V=\frac{(-1)^n}{2^n n !} \Omega^n\left(\text { resp. } d V=\frac{(-1)^n}{2^n n !} \eta \wedge(d \eta)^n\right).
     \end{equation}
\end{theorem}
\begin{remark}
    Not sure if it's correct up to some factors.
\end{remark}

\begin{definition}
    $\mathcal{A}$: the set of all associated metrics
    
    $\mathcal{N}$: the set of all metrics with the same volume element
    
    $\mathcal{M}$: the set of all metrics
\end{definition}
\begin{definition}
    For symmetric tensor fields $S$ and $T$ of type $(0,2)$ we define a Riemannian metric $(\cdot, \cdot)$ by
    \begin{equation}
        (S, T)_g=\int_M S_{i j} T_{k l} g^{i k} g^{j l} d V_g.
    \end{equation}
\end{definition}
\begin{remark}
    One can think of metrics with the same total volume on an n-dimensional manifold as being at a fixed distance from the zero tensor.
\end{remark}
\begin{proof}
    Consider the path $t g, t \in[0,1]$, from the zero tensor to the metric $g$; then since $g$ may be considered as the tangent to the path at each point, we have the following entertaining computation:
$$
|g| \equiv \int_0^1(g, g)_{t g}^{1 / 2} d t=\int_0^1\left(\frac{n}{t^2} \int_M \sqrt{t^n \operatorname{det} g} d x^1 \cdots d x^n\right)^{1 / 2} d t=4 \sqrt{\frac{\operatorname{vol}_g M}{n}}.
$$
\end{proof}


\begin{proposition}
    On a compact manifold $M$ the set $\mathcal{M}$ of all Riemannian metrics may be given a Riemannian metric (Ebin [1970]): The tangent space, $T_g \mathcal{M}$, of $\mathcal{M}$ at a metric $g$ is the space of symmetric tensor fields of type $(0,2)$.

    The tangent space $T_g\mathcal{A}$ of $\mathcal{A}$ at $g$ is the set of all symmetric tensor fields that anticommute with $J$ ($\phi$ and annihilate $\xi$).
\end{proposition}

\begin{theorem}
    Let $(M, \Omega)$ be a symplectic manifold, or respectively $(M, \eta)$ a contact manifold, and $f$ a diffeomorphism satisfying $f^* \Omega=\Omega$, respectively $f^* \eta=\eta$. Then for any associated metric $g, f^* g$ is also an associated metric.
\end{theorem}

\begin{lemma}
    Let $(M, g)$ be a compact orientable Riemannian manifold and for a vector field $V$ let $v(X)=g(V, X)$. Then $\int_M V^i \theta_i d V_g=0$ for every closed 1-form $\theta$ if and only if $v=\delta \Psi$ for some 2 -form $\Psi$. $\int_M v_i X^i d V_g=0$ for every vector field $X$ if and only if $v=0$.
\end{lemma}

\section{Examples of almost contact metric manifolds}

\begin{example}
    $\mathbb{R}^{2n+1}(-3).$ It's a contact metric structure on it.
\end{example}

\begin{example}
    $M^{2n+1}\subset \tilde{M}^{2n+2}$. Every $C^{\infty}$ orientable hypersurface of an almost complex manifold has an almost contact structure.

    In general, one cannot expect the induced almost contact metric structure to be a contact metric structure.
\end{example}
\begin{theorem}
    Let $M^{2 n+1}$ be a hypersurface of a Kähler manifold $\tilde{M}^{2 n+2},(\phi, \xi, \eta, g)$ its induced almost contact metric structure and $A$ its Weingarten map. Then $(\phi, \xi, \eta, g)$ is a contact metric structure if and only if 
    \begin{equation}
        A \phi+\phi A=-2 \phi.
    \end{equation}
\end{theorem}

\begin{example}
    $S^5\subset S^6$.
\end{example}

\begin{example}[The Boothby–Wang fibration]
   It has a \textbf{$K$-contact structure}.
\end{example}

\begin{definition}
    A contact metric structure for which the characteristic vector field is a Killing vector field is called a \textbf{$K$-contact structure}.
\end{definition}

\begin{theorem}
    Let $\pi: M^{2n+1}\to M^{2n}$ be the Boothby-Wang fibration of a compact regular contact manifold $M^{2n+1}$. Then the first Betti numbers of  $M^{2n+1}$ and $M^{2n}$ are equal.
\end{theorem}
\begin{theorem}
    No torus $T^{2n+1}$ can carry a regular contact structure.
\end{theorem}

\begin{example}
    $M^{2n}\times \mathbb{R}$. It is an almost contact metric structure that is certainly not a contact metric structure.
\end{example}

\begin{example}[Parallelizable manifolds]
In particular, any odd-dimensional Lie group carries an almost contact structure.
\end{example}
\begin{definition}
    A smooth manifold with or without boundary is said to be \textbf{parallelizable} if it admits a smooth global frame.
\end{definition}
\begin{theorem}
    A smooth vector bundle is smoothly trivial if and only if it admits a smooth global frame.
\end{theorem}


\section{Integral Submanifolds and Contact Transformations}

\subsection{Integral submanifolds}
Let $M^{2 n+1}$ be a contact manifold with contact form $\eta$. We have seen that $\eta=0$ defines a $2 n$-dimensional subbundle $\mathcal{D}$ called the \textbf{contact distribution} or\textbf{ subbundle} and that since $\eta \wedge(d \eta)^n \neq 0, \mathcal{D}$ is nonintegrable.

\begin{definition}
    A submanifold $M^r$ of $M^{2 n+1}$ is called an \textbf{integral submanifold} if $\eta(X)=0$ for every tangent vector $X$. It is clear that for any pair of tangent vector fields we have
$$
d \eta(X, Y)=\frac{1}{2}(X \eta(Y)-Y \eta(X)-\eta([X, Y])=0 .
$$

Then in terms of associated metrics, $g(X, \phi Y)=0$ and for this reason integral submanifolds are often called $C$-\textbf{totally real submanifolds.} In particular, $\phi$ maps tangent vectors to normal vectors; also, since $\xi$ is a normal vector, the dimension $r$ can be at most $n$. On the other hand, by the Darboux theorem, it's easy to see there exist an $n$-dim integral submanifold.
\end{definition}

\begin{theorem}
    Let $M^{2n+1}$ be a contact manifold with contact form $\eta$. Then there exist integral submanifolds of the contact subbundle $\mathcal{D}$ of dimension $n$ but of no higher dimension.
\end{theorem}

\begin{theorem}[Sasaki, 1964]
     Let $\left(x^i, y^i, z\right), i=1, \ldots, n$, be local coordinates about a point $m=\left(x_0^i,  y_0^i, z_0\right)$ such that $\eta=d z-\sum_{i=1}^n y^i d x^i$ on the coordinate neighborhood. In order that $r$ linearly independent vectors $X_\lambda, \lambda=$ $1, \ldots, r \leq n$, at $m$ with components $\left(a_\lambda^i, b_\lambda^i, c_\lambda\right)$ be tangent to an $r$ dimensional integral submanifold it is necessary and sufficient that $\eta\left(X_\lambda\right)=0$ and $d \eta\left(X_\lambda, Y_\mu\right)=0$, that is, $c_\lambda=\sum_i y_0^i a_\lambda^i$ and $\sum_i a_\lambda^i b_\mu^i=$ $\sum_i a_\mu^i b_\lambda^i$.
\end{theorem}

As in the symplectic case we note the abundance of integral submanifolds of $\mathcal{D}$.

\begin{theorem}[Sasaki, 1964]
    Given a vector $X \in \mathcal{D}$ at $m \in M^{2n+1}$ and any $r, 1 \leq r \leq n$, there exists an $r$-dimensional integral submanifold $M^r$ of $\mathcal{D}$ through m with $X$ tangent to $M^r$.
\end{theorem}

Similar to Weistein theorem, we have 
\begin{theorem}[Lychagin, 1977]
     If $L$ is an n-dimensional integral submanifold of a contact manifold $\left(M^{2 n+1}, \eta\right)$, then there exist an open neighborhood $\mathcal{U}$ of $L$ in $M^{2 n+1}$, an open neighborhood $\mathcal{V}$ of the zero section in $T^* L \times \mathbb{R}$, and a diffeomorphism $f: \mathcal{U} \longrightarrow \mathcal{V}$ such that $\left.f\right|_L$ is the identity on $L$ and $f^*(\beta-d t)=\eta$.
\end{theorem}

\subsection{Contact transformations}
Recall that a diffeomorphism $f$ of $M^{2 n+1}$ is a contact transformation if $f^* \eta=\tau \eta$ for some nonvanishing function $\tau$ and that $f$ is a strict contact transformation if $\tau \equiv 1$. Clearly $f^* d \eta=d \tau \wedge \eta+\tau d \eta$, so if $d \eta$ is invariant, $(\tau-1) d \eta=-d \tau \wedge \eta$ and hence $(\tau-1) \eta \wedge d \eta=0$, giving $\tau \equiv 1$. Thus $f$ is strict if and only if $d \eta$ is invariant.

\begin{theorem}
    $f$ is a contact transformation if and only if $X \in \mathcal{D}$ implies $f_* X \in \mathcal{D}$.
\end{theorem}

\begin{theorem}
    A diffeomorphism $f$ is a contact transformation if and only if $f$ maps $r$-dimensional integral submanifolds to $r$-dimensional integral submanifolds.
\end{theorem}

\begin{definition}
    If $L_{X}\eta=\sigma \eta$ for some function $\sigma$, $X$ is called an \textbf{infinitesimal contact transformation}.If $\sigma=0$,  we say that $X$ is a \textbf{strict infinitesimal contact transformation}.
\end{definition}

\begin{theorem}
    A vector field $X$ is an infinitesimal contact transformation if and only if there exists a function $f$ on $M^{2n+1}$ such that 
    \begin{equation}
        X = -\frac{1}{2} \phi \nabla f + f\xi.
    \end{equation}
\end{theorem}
\begin{proof}
    For sufficiency, we have 
    \begin{equation}
        L_{X}\eta=(\xi f) \eta.
    \end{equation}
    ...
\end{proof}

\begin{corollary}
    $X$ is strict if and only if $\xi f=0.$
\end{corollary}

As in the symplectic case we have the following theorem of Hatakeyama on the transitivity of the group of contact transformations.
\begin{theorem}[Hatakeyama, 1966]
    Let $ M^{2n+1}$ be a compact contact manifold. Then for any two points $p, q$, there exists a contact transformation mapping $p$ to $q$. If $M^{2n+1}$ is regular, there exists a strict contact transformation mapping $p$ to $q$.
\end{theorem}


\section{Sasakian and Cosymplectic Manifolds}

Let $M^{2 n+1}$ be an almost contact manifold with structure tensors $(\phi, \xi, \eta)$ and consider the manifold $M^{2 n+1} \times \mathbb{R}$. We denote a vector field on $M^{2 n+1} \times \mathbb{R}$ by $(X, f \frac{d}{d t})$, where $X$ is tangent to $M^{2 n+1}, t$ the coordinate on $\mathbb{R}$, and $f$ a $C^{\infty}$ function on $M^{2 n+1} \times \mathbb{R}$. Define an almost complex structure $J$ on $M^{2 n+1} \times \mathbb{R}$ by
$$
J\left(X, f \frac{d}{d t}\right)=\left(\phi X-f \xi, \eta(X) \frac{d}{d t}\right)
$$
that $J^2=-I$ is easy to check. If now $J$ is integrable, we say that the almost contact structure $(\phi, \xi, \eta)$ is \textbf{normal} (Sasaki and Hatakeyama [1961]).

\begin{remark}
    Compare to the ``regular''.
\end{remark}

\begin{definition}
    We are thus led to define four tensors $N^{(1)}, N^{(2)}, N^{(3)}, N^{(4)}$ by
    \begin{equation}
        \begin{aligned}
N^{(1)}(X, Y) & =[\phi, \phi](X, Y)+2 d \eta(X, Y) \xi, \\
N^{(2)}(X, Y) & =\left(L_{\phi X} \eta\right)(Y)-\left(L_{\phi Y} \eta\right)(X), \\
N^{(3)} & =\left(L_{\xi} \phi\right) X \\
N^{(4)} & =\left(L_{\xi} \eta\right)(X) .
\end{aligned}
    \end{equation}
\end{definition}

\begin{theorem}
    For an almost contact structure $(\phi, \xi, \eta)$ the vanishing of $N^{(1)}$ implies the vanishing of $N^{(2)}, N^{(3)}$ and $N^{(4)}$.
\end{theorem}
\begin{theorem}
     For a contact structure $(\phi, \xi, \eta, g)$, $N^{(2)}$ and $N^{(4)}$ vanish. Moreover, $N^{(3)}$ vanishes if and only if $\xi$ is a Killing vector field.
\end{theorem}

\begin{lemma}
     For an almost contact metric structure $(\phi, \xi, \eta, g)$, the covariant derivative of $\phi$ is given by
     \begin{equation}
         \begin{aligned}
2 g\left(\left(\nabla_X \phi\right) Y, Z\right)= & 3 d \Phi(X, \phi Y, \phi Z)-3 d \Phi(X, Y, Z)+g\left(N^{(1)}(Y, Z), \phi X\right) \\
& +N^{(2)}(Y, Z) \eta(X)+2 d \eta(\phi Y, X) \eta(Z) \\
& -2 d \eta(\phi Z, X) \eta(Y).
\end{aligned}
     \end{equation}
\end{lemma}

\begin{definition}
Define  a tensor
    \begin{equation}
        h:=\frac{1}{2}\left(L_{\xi} \phi\right) X=\frac{1}{2}N^{(3)}.
    \end{equation}
\end{definition}
\begin{remark}
    The first property to note is immediate, namely $h \xi=0$. 
\end{remark}

\begin{lemma}
    On a contact metric manifold $h$ is a symmetric operator,
    \begin{equation}
        \nabla_X \xi=-\phi X-\phi h X.
    \end{equation}
$h$ anticommutes with $\phi$, and $\operatorname{tr} h=0$.
\end{lemma}
\begin{corollary}
    On a contact metric manifold $\delta \eta=0$.
\end{corollary}

\subsection{Definition of Sasakian manifolds}

\begin{definition}
    A Sasakian manifold is a normal contact metric manifold.
\end{definition}
\begin{theorem}
    An almost contact metric structure $(\phi, \xi, \eta, g)$ is Sasakian if and only if
    \begin{equation}\label{eq28}
        \left(\nabla_X \phi\right) Y=g(X, Y) \xi-\eta(Y) X .
    \end{equation}
\end{theorem}

\begin{corollary}
    A Sasakian manifold is $K$-contact.
\end{corollary}
\begin{remark}
    In dimension 3 the converse is true, Corollary 6.5. For K-contact structures that are not Sasakian, see Example 6.7.2.
\end{remark}

Let's draw two diagrams to make a conclusion:

\begin{tikzpicture}  
\node[block] (a) {Hermitian};   
\node[block,right=80pt of a] (b) {Kaehler};  
 % the commands used for the different location of different blocks   
 \node[block] (c) at ([yshift=-3cm]$(a)!1.0!(a)$) {almost Hermitian};  
  \node[block] (d) at ([yshift=-3cm]$(b)!1.0!(b)$) {almost Kaehler};   
  \draw[line] (a)--node[above]{$d\Omega =0$} (b);  
  \draw[line] (d)-- node[left]{$[J,J]=0$} (b);  
  \draw[line] (c)-- node[left]{$[J,J]=0$} (a);   
  \draw[line] (c)-- node[above]{$d\Omega =0$}(d);
  \draw[line] (c)--node[left]{$\nabla J=0$}  (b);
\end{tikzpicture}  
\begin{tikzpicture}  
\node[block] (a) {normal almost contact metric};   
\node[block,right=80pt of a] (b) {Sasakian};  
 % the commands used for the different location of different blocks   
 \node[block] (c) at ([yshift=-3cm]$(a)!1.0!(a)$) {almost contact metric};  
  \node[block] (d) at ([yshift=-3cm]$(b)!1.0!(b)$) { contact metric};   
  \draw[line] (a)--node[above]{$\Phi=d \eta$} (b);  
  \draw[line] (d)-- node[left]{$N^{(1)}=0$} (b);  
  \draw[line] (c)-- node[left]{$N^{(1)}=0$} (a);   
  \draw[line] (c)-- node[above]{$\Phi=d \eta$}(d);
  \draw[line] (c)--node[left]{$\eqref{eq28}$}  (b);
\end{tikzpicture}  





\subsection{More about Sasakian manifolds}

Recall that
    a \textbf{contact metric structure} on an odd-dimensional manifold $N^{2 n+1}$ is given by $(\varphi, \xi, \eta, g)$, where $\varphi$ is a tensor field of type $(1,1)$ on $N, \xi$ is a vector field, $\eta$ is an 1-form and $g$ is a Riemannian metric such that
    \begin{equation}\label{eq1}
        \varphi^2=-I+\eta \otimes \xi, \quad \eta(\xi)=1
    \end{equation}
and
\begin{equation}\label{eq2}
    g(\varphi X, \varphi Y)=g(X, Y)-\eta(X) \eta(Y), \quad g(X, \varphi Y)=d \eta(X, Y), \quad \forall X, Y \in C(T N) .
\end{equation}

\begin{remark} Recall that \cite{dragomir_differential_2006}, from $\eqref{eq1}$, one can deduce
    $$
\varphi \xi=0, \quad \eta \circ \varphi=0.
$$
\end{remark}

\begin{example}
    Strictly pseudoconvex CR manifolds are therefore contact Riemannian manifolds, in a natural way.
\end{example}
\begin{theorem}
   Let $\left(M, T_{1,0}(M), \theta\right)$ be a pseudo-Hermitian manifold. Then the Webster metric $g_\theta$ is Sasakian if and only if the Tanaka-Webster connection of $\left(M, T_{1,0}(M), \theta\right)$ has vanishing pseudo-Hermitian torsion, that is, $\tau=0$.
\end{theorem}

Recall that
    a contact metric manifold $(N, \varphi, \xi, \eta, g)$ is called \textbf{Sasakian} if it is normal, i.e.
$$
N_{\varphi}+2 d \eta \otimes \xi=0,
$$
where
\begin{equation}\label{eq:Jtorsion}
    N_{\varphi}(X, Y)=[\varphi X, \varphi Y]-\varphi[\varphi X, Y]-\varphi[X, \varphi Y]+\varphi^2[X, Y], \quad \forall X, Y \in C(T N),
\end{equation}
is the \textbf{Nijenhuis tensor} field of $\varphi$, or, equivalently, if
\begin{equation}
    \left(\nabla_X \varphi\right)(Y)=g(X, Y) \xi-\eta(Y) X, \quad \forall X, Y \in C(T N) .
\end{equation}

\begin{remark}
    We note that from the above formula it follows 
    \begin{equation}
        \nabla_X \xi=-\varphi X.
    \end{equation}
\end{remark}

\subsection{Connections and curvatures}

\begin{theorem}
 Let $\left(M, T_{1,0}(M), \theta\right)$ be a pseudo-Hermitian manifold and $g_\theta$ the Webster metric of $\left(M, T_{1,0}(M), \theta\right)$. Then there exists a unique linear connection $\nabla$ on $M$, called the \textbf{Tanaka-Webster} connection, such that:
 
(1) the Levi distribution $H(M)$ is parallel with respect to $\nabla$;

(2) $\nabla g_\theta=0, \nabla J=0, \nabla \theta=0$ (hence $\left.\nabla T=0\right)$;

(3) the torsion $T_{\nabla}$ of $\nabla$ satisfies $T_{\nabla}(X, Y)=-2 \Omega(X, Y) T$ and $T_{\nabla}(T, J X)=-J T_{\nabla}(T, X)$, for any $X, Y \in H(M)$.
\end{theorem}

We denote by $\nabla^\theta$ the Levi-Civita connection of $\left(M, g_\theta\right)$. One can see that the Levi-Civita connection $\nabla^\theta$ is related to the Tanaka-Webster connection $\nabla$ by
\begin{equation}\label{eq33}
    \nabla^\theta=\nabla+(\Omega-A) \otimes T+\tau \otimes \theta+2 \theta \odot J,
\end{equation}
where $2(\theta \odot J)(X, Y)=\theta(X) J Y+\theta(Y) J X$. Thus, we have $\nabla_X^\theta T=$ $\tau(X)+J X$. In particular, 
\begin{equation}
    \nabla_T^\theta T=0.
\end{equation}
If $X, Y \in H(M)$, then
\begin{equation}
    \nabla_X^\theta Y=\nabla_X Y+[\Omega(X, Y)-A(X, Y)] T .
\end{equation}
Since $\nabla^\theta-\nabla$ is a $(1,2)$ tensor field on $M$ and $\nabla_T^\theta T=\nabla_T T=0$, we can define a vector field $V$ on $M$ given by
\begin{equation}
    V=\operatorname{trace}_{g_\theta}\left(\nabla^\theta-\nabla\right)=\operatorname{trace}_{G_\theta}\left(\nabla^\theta-\nabla\right) .
\end{equation}
Actually we have 
\begin{equation}
    V=-\operatorname{trace}_{g_\theta}(\tau) T=0
\end{equation}







\begin{definition}
    Let $(N, \varphi, \xi, \eta, g)$ be a Sasakian manifold. The sectional curvature of a 2 -plane generated by $X$ and $\varphi X$, where $X$ is an unit vector orthogonal to $\xi$, is called the $\varphi$-\textbf{sectional curvature} determined by $X$. If the $\varphi$-sectional curvature is a constant $c$, then $(N, \varphi, \xi, \eta, g)$ is called a \textbf{Sasakian space form} and it is denoted by $N(c)$.
The curvature tensor field of a Sasakian space form $N(c)$ is given by
\begin{equation}
    \begin{aligned}
R(X, Y) Z= & \frac{c+3}{4}\{g(Z, Y) X-g(Z, X) Y\}+\frac{c-1}{4}\{\eta(Z) \eta(X) Y- \\
& -\eta(Z) \eta(Y) X+g(Z, X) \eta(Y) \xi-g(Z, Y) \eta(X) \xi+ \\
& +g(Z, \varphi Y) \varphi X-g(Z, \varphi X) \varphi Y+2 g(X, \varphi Y) \varphi Z\}
\end{aligned}
\end{equation}
\end{definition}

\begin{remark}[cf. \cite{chen_bochner-type_2012}]
    The $\varphi$-sectional curvature determines the curvature tensor $R$ completely.
\end{remark}

\begin{remark}
    A direct consequence of $\eqref{eq33}$ is that, for Sasakian manifolds, we have(substitute $\tau=0, A=0$ )
    \begin{equation}
        K(X, JX)=K^{\theta}(X,JX)+3, \forall X \in \mathcal{H}(M),
    \end{equation}
    and 
    \begin{equation}
        K(X, T)=K^{\theta}(X,T)-1=0, \forall X \in \mathcal{H}(M),
    \end{equation}
    where $K$ is sectional curvature wrt TW connection, while $K^{\theta}$ is the sectional curvature wrt LC connection.
\end{remark}

\begin{theorem}[cf. \cite{blair_riemannian_2010}]
     Let $M(c)$ be a complete, simply connected Sasakian manifold with constant $\phi$-sectional curvature $c$. Then $M(c)$ belongs one of the three families of examples listed below.
\end{theorem}

To begin, let $(\phi, \xi, \eta, g)$ be a contact metric structure and recall the notion of a $\mathcal{D}$-\textbf{homothetic deformation}:
\begin{equation}\label{eq6}
    \bar{\eta}=a \eta, \quad \bar{\xi}=\frac{1}{a} \xi, \quad \bar{\phi}=\phi, \quad \bar{g}=a g+a(a-1) \eta \otimes \eta,
\end{equation}
where $a$ is a positive constant. The deformed structure $(\bar{\phi}, \bar{\xi}, \bar{\eta}, \bar{g})$ is again a contact metric structure, and it enjoys many of the properties of the original structure. In particular, if $(\phi, \xi, \eta, g)$ is Sasakian, so is $(\bar{\phi}, \bar{\xi}, \bar{\eta}, \bar{g})$; if $M(c)$ is a Sasakian space form, then deforming the structure, we obtain the Sasakian space form $M(\bar{c})$ where $$\bar{c}=\frac{c+3}{a}-3$$ (see Tanno [1968], [1969] for details). We will show that there exist Sasakian space forms $M(c)$ for every value of $c$. 

\begin{remark}
    $\eqref{eq6}$ is natural from my point of view: we want to make a conformal transformation on $\mathcal{D}$ or ,more precisely, homothetic deformation,  so we hope when restricting to $\mathcal{D}$: $\bar{g}=ag$, so we may assume $\bar{g}=ag+b\eta\otimes \eta.$ But from the 2nd of $\eqref{eq2}$, we are suggested to let $\bar{\eta}=a\eta$ so $\bar{\xi}$ follows its step. Now put the above transformations to the 1st of $\eqref{eq2}$ (testing $\bar{\xi}$), we get $b=a^2-a$.
\end{remark}

\begin{example}[$\mathbb{S}^{2 n+1}, c>-3$]
    Let $\mathbb{S}^{2 n+1}=\left\{z \in \mathbb{C}^{n+1}:|z|=1\right\}$ be the unit $(2 n+1)$-dimensional Euclidean sphere. Consider the following structure tensor fields on $\mathbb{S}^{2 n+1}$ : the standard metric field $g_0$, the vector field $\xi_0(z)=-\mathcal{J} z, z \in \mathbb{S}^{2 n+1}$, where $\mathcal{J}$ is the usual almost complex structure on $\mathbb{C}^{n+1}$ defined by
$$
\mathcal{J} z=\left(-y^1, \ldots,-y^{n+1}, x^1, \ldots, x^{n+1}\right),
$$
for $z=\left(x^1, \ldots, x^{n+1}, y^1, \ldots, y^{n+1}\right)$, and $\varphi_0=s \circ \mathcal{J}$, where $s: T_z \mathbb{C}^{n+1} \rightarrow T_z \mathbb{S}^{2 n+1}$ denotes the orthogonal projection.  Then $g_0$ and $\xi_0$ detemine $\eta_0$ and $\varphi_0$ by 
$$
\eta_0=g_0(\xi_0,\cdot), \quad d\eta_0(X,Y)=2g( X,\varphi_0 Y).
$$
Equipped with these tensors, $\mathbb{S}^{2 n+1}$ becomes a Sasakian space form with the $\varphi_0$-sectional curvature equal to 1 , denoted by $\mathbb{S}^{2 n+1}(1)$. Now, consider the deformed structure on $\mathbb{S}^{2 n+1}$
$$
\eta=a \eta_0, \quad \xi=\frac{1}{a} \xi_0, \quad \varphi=\varphi_0, \quad g=a g_0+a(a-1) \eta_0 \otimes \eta_0
$$
where $a$ is a positive constant. The structure $(\varphi, \xi, \eta, g)$ is still a Sasakian structure and $\left(\mathbb{S}^{2 n+1}, \varphi, \xi, \eta, g\right)$ is a Sasakian space form with constant $\varphi$-sectional curvature $c=\frac{4}{a}-3$, $c>-3$, denoted by $\mathbb{S}^{2 n+1}(c)$.
\end{example}

\begin{example}[$\mathbb{R}^{2 n+1}, c=-3$]
    The Euclidean space $\mathbb{R}^{2 n+1}\left(x^1, \ldots, x^n, y^1, \ldots, y^n, z\right)$ is a Sasaki manifold with Sasaki structure $(\xi, \eta, \phi, g)$ as follows:
$$
\begin{gathered}
\xi=2 \frac{\partial}{\partial z}, \quad \eta=\frac{1}{2}\left(d z-\sum_{i=1}^n y^i d x^i\right), \quad g=\eta \otimes \eta+\frac{1}{4} \sum_{i=1}^n\left(\left(d x^i\right)^2+\left(d y^i\right)^2\right), \\
\phi\left(\frac{\partial}{\partial x^i}\right)=-\frac{\partial}{\partial y^i}, \quad \phi\left(\frac{\partial}{\partial y^i}\right)=\frac{\partial}{\partial x^i}+y^i \frac{\partial}{\partial z}, \quad \phi\left(\frac{\partial}{\partial z}\right)=0 .
\end{gathered}
$$
\end{example}

\begin{example}[$B^n \times \mathbb{R}, c<-3$]
    Let $B^n$ be a simply connected bounded domain in $\mathbb{C}^n$ and $(J, G)$ a Kähler structure with constant holomorphic sectional curvature $k<0$. For such a structure, the fundamental 2 -form $\Omega$ of the Kähler structure is exact and hence $\Omega=d \omega$ for some real analytic 1 -form $\omega$. Now on $B^n \times \mathbb{R}$ let $\pi$ denote the projection onto $B^n$ and $t$ the coordinate on $\mathbb{R}$. Then $B^n \times \mathbb{R}$ with $\eta=\pi^* \omega+d t$ and $g=\pi^* G+\eta \otimes \eta$ is a Sasakian manifold. Regarding $\eta$ as a connection form on $B^n \times \mathbb{R}$, let $\tilde{\pi} X$ denote the horizontal lift of a vector field $X$ on $B^n$. Also we denote by $\underline{K}$ the sectional curvature of $B^n$. Then by direct computation (see Ogiue [1965] or use the Riemannian submersion technique of O'Neill [1966]), we have $K(\tilde{\pi} X, \tilde{\pi} Y)=\underline{K}(X, Y)-3 \eta\left(\nabla_{\tilde{\pi} X} \tilde{\pi} Y\right)^2$ where $\{X, Y\}$ is an orthonormal pair on $B^n$. Now $g\left(\nabla_{\tilde{\pi} X} \tilde{\pi} Y, \xi\right)=-g\left(\tilde{\pi} Y, \nabla_{\tilde{\pi} X} \xi\right)=g(\tilde{\pi} Y, \phi \tilde{\pi} X)=$ $g(\tilde{\pi} Y, \tilde{\pi} J X)$. Therefore $B^n \times \mathbb{R}$ has constant $\phi$-sectional curvature $c=$ $k-3$
\end{example}

The following is a Sasakian analogue of Schur’s Theorem.
\begin{theorem}[cf. \cite{adachi_geometric_2010}]
     If the $\phi$-sectional curvature at each point of a Sasakian manifold $M$ of dimension $\geq 5$ does not depend on the choice of $\phi$-section at that point, then it is constant on $M$.
\end{theorem}

The following is the unique
existence theorem of Sasakian space forms.
\begin{theorem}
    For any two simply connected complete Sasakian manifolds of constant $\phi$-sectional curvature $c$, there exists an isomorphism between them which preserves their almost contact metric structures.
\end{theorem}

